\documentclass[11pt]{article}
\usepackage[utf8]{inputenc}
\usepackage[T1]{fontenc}
\usepackage{lmodern}
\usepackage[margin=1in]{geometry}
\usepackage{csquotes}
\usepackage[style=authoryear,backend=biber,sorting=nyt,maxbibnames=99,giveninits=true,uniquename=init]{biblatex}
\usepackage[hidelinks]{hyperref}
\addbibresource{references.bib}
\begin{document}
\section*{Subsurface Drainage for Soil Salt Leaching}

\subsection{Overview}

Subsurface (tile) drainage is a cornerstone technology for reclaiming and maintaining salt-affected agricultural land, yet the literature converges on a striking recognition that the salt-leaching outcomes it produces are only partially understood and, in many settings, only partially detectable with current observational programs. The recurring message across studies is temporal: meaningful salt reduction often accrues slowly, so that short experimental windows may underestimate long-term benefits, and only extended drainage periods combined with sustained monitoring can resolve the true trajectory of desalinization. \textcite{Almen2021} state plainly that additional years of extensive drainage may be needed before noticeable field-scale salt reduction appears, while \textcite{Daigh2025} argue that longer-term perspectives should be incorporated when assessing soil-management efficacy, since immediate gains may understate protective benefits against salinization. This frontier is therefore defined less by settled magnitudes than by open questions concerning (i) the temporal and mechanistic understanding of salt removal, (ii) the instrumentation and monitoring infrastructure required to close water and salt balances, (iii) sensing, mapping, and modelling tools for quantifying and forecasting salinity, (iv) design and management optimization across heterogeneous soils and climates, (v) the fate, treatment, and safe disposal of saline and contaminant-laden drainage water, and (vi) the scaling of understanding from plots to landscapes and watersheds. The evidence base spans arid and coastal reclamation zones in China, the Indian subcontinent, the Mediterranean and Middle East, Australia, northern Europe, and the drainage-intensive and evaporation-pond regions of North America, and the open questions differ in emphasis but share a common demand for longer records, better process representation, and more transferable relationships.

\subsection{Long-term monitoring and the temporal dynamics of desalinization}

A dominant research gap is the need to characterize salt removal over multi-year horizons rather than single seasons. Beyond the call by \textcite{Almen2021} for additional years of drainage before field-scale salt reduction becomes evident, several authors emphasize that the sustainability of combined saline-water irrigation and subsurface drainage systems can only be judged over the long term: \textcite{Feng2017} argue that the long-term effects of saline-water irrigation and subsurface drainage on soil salinity and crop yield require further evaluation, and \textcite{Feng2021} similarly urge that this combination receive more attention for long-term soil-salinity and grain-yield management. Desalination rates and their trends themselves remain incompletely quantified; \textcite{Guo2024b} identify both the rates and their temporal trends as requiring additional study, and further note that the suitability of such systems in arid areas with deeper groundwater levels needs exploration. \textcite{Yao2024} call for investigation of the long-term effects of specific drainage parameters on agricultural output and soil health. Controlled drainage introduces its own unresolved temporal trade-offs for salinity: \textcite{Jouni2018} state repeatedly that more research is needed on the long-term effects of controlled drainage on soil salinity and on optimizing management options to avoid salt accumulation. Soil-degradation processes under drainage likewise demand extended experiments, as \textcite{Lu2019} note, coupled with site-specific cost--benefit assessments under particular hydrological and salinity conditions. In coastal, heavy-textured saline soils, \textcite{Islam2025} identify the absence of long-term studies on nutrient loss and water pollution associated with surface and subsurface drainage as a specific deficiency. \textcite{Daigh2025} frame this collective need as a matter of landscape resilience, advocating continued case studies of how soil-management practices influence resilience to salinization and cautioning against underestimating protective benefits from short records. At the salt-balance level, \textcite{Heng2021} recommend regular sampling surveys to track intensifying salt accumulation and the formulation of salt-leaching plans keyed to salt-balance trends and salt-leaching curves.

\subsection{Instrumentation, surface-water monitoring, and closing salt balances}

Several studies argue that the field cannot advance without better measurement of what leaves the system through the drains. \textcite{Mastrocicco2013} contend that much greater attention must be paid to surface-water monitoring to assess the potential export of salts and fertilizers through subsurface drainage. \textcite{Askri2022} make the instrumentation requirement concrete for oasis-scale water and salt balances, calling for the installation of flow meters and electrical-conductivity probes to measure drainage discharge and salt concentration continuously and automatically; they further note that quantifying these balance components for recent years would reveal the long-term performance of the drainage system, that the influence of the system on subsurface water and salt dynamics is still poorly understood, and that the impact of saline drainage water on downstream phreatic groundwater quality must be evaluated to safeguard sustainability. The environmental accounting of drainage improvement is also unresolved in nutrient terms: \textcite{Turtola1995} argue that the environmental benefits or disadvantages of subsurface drainage improvement should be evaluated through the nutrient balances of recipient waters. Monitoring infrastructure for tracking salinity trends over time is another recurring recommendation; \textcite{Singh2020d} propose installing representative cluster wells within experimental fields to track salinity at two depths, and \textcite{Gafni2001} plan deep drilling and soil sampling to monitor salinity dynamics over time.

\subsection{Process understanding of salt leaching}

\subsubsection{Preferential flow, macropores, and heterogeneous media}

The physical mechanisms controlling how salts are mobilized and removed remain an active frontier. \textcite{Xin2016} highlight two open questions: how macropores affect salt leaching under transient-state conditions, and how leaching methods might be modified to avoid or minimize macropore effects. For layered or otherwise structured profiles, \textcite{Zhang2022} call for future studies to incorporate the governing factors of salt leaching in low-permeability or structured soil media. \textcite{Li2023} argue that both field experiments and numerical simulations are needed to investigate the soil-salt replacement effect between subsurface pipes under brackish-water leaching and to reconcile it with indoor simulation results.

\subsubsection{Rainfall, irrigation, and desalinization ratios}

The interaction of natural and applied water with drainage is central to leaching efficiency. \textcite{Yu2015} suggest that greater rainfall frequency, irrigation, and subsurface drainage sustained over several years may improve the homogenizing effect on soil salinity, and that increasing the intensity of rainfall and subsurface drainage could raise the desalinization ratio. \textcite{Chen2025} propose analyzing variable scenarios to derive more generally applicable relationships between the soil-desalinization process and rainfall/irrigation requirements under different drainage conditions, and extending such work to different climatic conditions and soil types. \textcite{Wang2024d} emphasize that evaporative water loss during summer waterlogging leaching must be fully accounted for when adjusting leaching-water amounts to achieve target results.

\subsubsection{Ponding strategy and leaching curves}

Where reclamation proceeds by ponding, \textcite{RAO1991} propose specific untested strategies: complete ponding followed by ponding over decreasing surface fractions (80\%, 60\%, 40\%, 20\%) to avoid unleached zones, and, conversely, reversing the sequence by starting at 20\% and increasing to complete ponding to counter re-salinization and economize water; they also suggest crop-residue mulches or cultivation techniques to control capillary rise in partially ponded situations, framing this as a rationale for future leaching experiments in drained saline soils of Haryana, India. Complementing this, \textcite{Yang2022c} stress the importance of determining soil-salinity leaching curves and generating fitting equations across different agricultural regions.

\subsubsection{Rootzone salinity, matric potential, and crop response}

Linking leaching to plant performance remains only loosely resolved. \textcite{Zhang2021} call for research to optimize irrigation depths for a given irrigation-water salinity, to quantify the effects of soil matric potential and water salinity on soil salinity and plant growth, and to determine the optimal rootzone matric-potential threshold; the authors also seek comparison of gravel-sand-layer drainage against conventional tile systems. \textcite{Heng2021} identify the link between cotton yield and soil nutrients under saline-alkali improvement, and the use of moderate deficit irrigation with subsurface pipe drainage (MDI-SPD) to reduce salinity with minimal water while raising yield, as open questions.

\subsection{Sensing, mapping, and modelling salinity}

A large cluster of frontier work concerns tools to observe and predict salinity at scale. Electromagnetic induction (EMI) methods are seen as promising but incompletely developed: \textcite{Narjary2024} propose using EMI for characterizing and mapping soil salinity, extending it to three-dimensional representation via inversion modelling (EM4soil), while cautioning that calibrated relationships must be validated and modified for other soil orders and diverse agro-climatic and hydrogeological settings. \textcite{Yu2025} argue that coupled apparent-electrical-conductivity--water-content (ECa--$\theta$) models could better estimate soil-solution salinity and refine the quantification of leaching rates and their mechanisms. \textcite{Singh2020d} envisage continuation of EM38 mapping programs through university collaborations and funding, and the use of multispectral, drone, and hyperspectral remote sensing to improve salinity assessment. Modelling frameworks require targeted extension: \textcite{Singh2020d} call for a new CSUID-1D model interface incorporating tile drains, drain depth, and drainage yield, while \textcite{Jayawardane2001} report modelling of water and salt flow through FILTER plots to explain progressive changes in soil and drainage-water salinity, and note that modelling sodicity and salinity changes can predict longer-term soil stability and the need for remedial measures such as periodic gypsum application. The most integrative vision is offered by \textcite{Zhu2025b}, who advocate developing zonal drainage-design algorithms based on the spatial variability of soil texture, quantifying the contribution of crop root-zone hydraulic enhancements to salinity redistribution, and constructing a multi-scale water-and-salt transport model coupled with remote-sensing inversion data for real-time salinity early warning and dynamic control; they add that soil-texture mapping and groundwater-salinity thresholds will remain necessary to guide precise drainage zoning while optimizing crop--drainage synergies through root-dynamic modelling.

\subsection{Design optimization and management under heterogeneity}

Translating process understanding into practice requires site- and texture-specific design that is not yet standardized. \textcite{vandenEertwegh2006} argue that changes in subsurface-drainage system redesign, surface-water-level management, and field-scale nutrient management must be integrated rather than treated in isolation. \textcite{Neha2022} call for more precise understanding of water and salt dynamics and the economic viability of cut-soiler preferential subsurface drainage (PSSD) across different textural classes of salt-affected soils, variable irrigation-water salinity, and different crop-residue types. \textcite{Wang2024d} suggest that appropriately increasing vertical porous flow pipes filled with straw or sand can enhance the salt-drainage effect of subsurface pipes, and \textcite{Yao2024} note that the effects of fully charged pipe segments on salt leaching require further investigation. \textcite{Jayawardane2001} caution that the optimum combination of practices will depend on specific site conditions, echoing the broader recognition that transferable design rules remain elusive.

\subsection{Fate, treatment, and disposal of saline and contaminant-laden drainage water}

Removing salts from the soil relocates the problem to the drainage water, and its management is among the most persistent open challenges. \textcite{Ayars2014} state simply that research is needed to find environmentally acceptable disposal methods, and \textcite{Steppuhn2006} note that additional disposal techniques may be necessary when drainage discharge volumes or dissolved-salt concentrations are large. Evaporation-pond management is a specific frontier: \textcite{BRADFORD1991} call for testing of pond-management methods under laboratory and field conditions representative of full-scale evaporation ponds, and suggest combining pond management with measures to reduce drainage volume or contamination---water conservation, salt-tolerant and drought-resistant crops, cycling water through salt-resistant vegetation, and treatment to remove contaminants. Trace-element fate compounds the disposal problem: \textcite{Gao2007} argue that the role of sediment mechanisms in the fate of arsenic in drainage disposal requires further examination of As speciation, transformation, and partitioning in sediments, while \textcite{Duff1997} propose investigating the feasibility of uranium recovery, on-site harvesting of pond algae to lower water and sediment U concentrations, and removal of the surface 10~cm of pond sediment to significantly reduce pond U levels. Emerging treatment routes are only at the demonstration stage: \textcite{Yang2023} outline a research program to verify the electrolysis characteristics of saline wastewater from subsurface pipes at different salinity levels and compositions, to perform wastewater-electrolysis experiments, to evaluate solar-energy resources for agricultural wastewater treatment, and to establish demonstration bases for large-scale application. \textcite{Pugliese2023} call for further investigation to enhance the sustainability and cost-effectiveness of sorptive filter materials---which can be recycled as soil amendment or regenerated---and stress that accurate prediction of filter performance requires testing with real drainage water rather than artificial drainage water. Blending strategies remain a longer-term goal, with \textcite{Singh2020d} aiming to customize the blend of subsurface drainage water and reverse-osmosis product water to support sustainable forage production.

\subsection{Scaling from plots to landscapes and watersheds}

A final frontier concerns the extrapolation of field observations to the regional, landscape, and watershed scales at which salinity is ultimately managed. \textcite{Yu2025} argue that future research should prioritize regional-scale planning of agri-aquaculture systems to avoid ineffective management, and that landscape-level water and salinity management will be crucial for informing best practices in coastal reclamation and other salinity-sensitive land uses. \textcite{Chen2025} frame a broader ecological dimension, proposing exploration of the ecological benefits gained from early cropping production on saline soils. At the watershed scale, \textcite{Frisbee2022} call for repeat synoptic sampling campaigns to characterize baseflow during a normal precipitation year and to identify the contribution of deep-aquifer groundwater to baseflow at downstream sites, for additional campaigns combined with groundwater-well sampling to understand seasonal variability in baseflow and salinity sources, for coupling $^{36}$Cl/Cl with other tracers to disentangle the impacts of road salt versus KCl fertilizers, and for further research on the effect of KCl fertilizers on the overall salinity budgets of agricultural watersheds. Together with the zonal, multi-scale modelling agenda of \textcite{Zhu2025b} and the arid-area suitability questions raised by \textcite{Guo2024b}, these studies define a clear trajectory: from characterizing salt removal in individual drained fields toward predicting and managing salinity across heterogeneous soils, climates, and connected hydrological landscapes.

\printbibliography
\end{document}
